Schuld, schaamte en schande zijn de kroonjuwelen van het Boeddhisme

Het kan tegenwoordig vreemd lijken om schaamte te voelen over je interesse in of gevoelens rondom seks. Als moderne, verlichte mensen zouden we allemaal extreem zelfverzekerd, goed aangepast en enthousiast moeten zijn over het onderwerp seks.
Maar niets is minder waar. Seksuele schaamte is, in werkelijkheid, voor velen van ons helemaal niet verdwenen. Het is vooral een psychologisch probleem en niet zomaar een politiek of religieus vraagstuk.
Onze capaciteit om ons seksueel op een zelfverzekerde en gelukkige manier uit te drukken, onze mogelijkheid om te zeggen wat we willen, er zonder schaamte om te vragen en snel situaties te verlaten waarin we ons niet vervuld of vernederd voelen, zijn allemaal enorme psychologische prestaties. Deze zaken zijn vaak alleen spontaan toegankelijk voor degenen die een zeer ondersteunende en emotioneel ontwikkelde vroege omgeving hebben ervaren. Om van nature onbezwaarde, seksueel gelukkige volwassenen te zijn, is nodig dat anderen, die zelf ontspannen waren, ons lang geleden hebben laten voelen dat we acceptabel zijn zoals we zijn. Dat onze lichamen en hun functies natuurlijke en goede dingen zijn. Dat we niet stout of zondig waren omdat we nieuwsgierigheid uitten over lichamelijk genot. Dat het oké was om af en toe rommel te maken - en dat het bijvoorbeeld meer dan goed was om op de leeftijd van twee jaar echt gefascineerd en verrukt te zijn door het vreemde en wonderlijke bestaan van je eigen achterwerk.
Seksuele verlangens behoren tot de meest persoonlijke en kwetsbare dingen die we ooit worden uitgenodigd om te uiten, en ze stellen ons bloot aan mogelijke enorme niveaus van spot en schaamte. Zoals pestkoppen van alle soorten altijd hebben geweten: als je iemand snel wilt vernietigen, maak hem dan beschaamd over zijn seksualiteit; hij zal nooit meer het zelfvertrouwen hebben om je uit te dagen.
Er zijn weinig dingen die meer wezenlijk ‘onszelf’ zijn dan ons verlangen naar seksuele verbinding, en daarom komen onzekerheden over hoe aardig we zijn, hoe verdienend we zijn of hoe legitiem het is dat we bestaan, vaak sterk naar voren in de slaapkamer. Ze vernietigen ons vermogen om eenvoudige en onbewogen seksuele wezens te zijn. Om het simpel te stellen: als er enig risico bestaat dat we ons slecht voelen over onszelf, voelen we ons - door een psychologische onvermijdelijkheid - slecht over onszelf en seks. Wat seksuele problemen worden genoemd - zoals impotentie, vaginisme, een gebrek aan verlangen, schadelijke verslavingen of een algemene angst voor intimiteit - zijn eerst en vooral problemen van zelfhaat. In de regel kun je niet tegelijkertijd jezelf haten en een geweldige tijd in bed hebben.
Het begin van het oplossen van het probleem van seksuele schaamte berust op een basisacceptatie dat het probleem bestaat en waarschijnlijk al veel schade heeft aangericht in ons leven. We moeten leren het te benoemen en te erkennen, tegen onszelf en een paar dierbaren te zeggen: ‘Ik voel verlammende schaamte rondom seks, en dat is oké.’ Een toewijding aan verandering is wat telt. Ondanks alle vrolijke suggesties die het tegendeel beweren, lopen veel van ons, vrouwen en mannen, op dit moment (net als tijdens het hoogtepunt van de Spaanse Inquisitie) rond met intense seksuele schaamte. Dit is niet omdat wat we seksueel willen op objectieve manier ‘slecht’ is (dat wil zeggen, opzettelijk schadelijk voor iemand anders), maar omdat onze geschiedenis ons heeft gepredisponeerd om zo negatief te voelen over onszelf.
Een belangrijk effect van seksuele schaamte is dat het ons tot zwijgen brengt. We zijn zo beschaamd dat we zelfs onze schaamte niet kunnen uitspreken. Het is daarom van groot belang om het lef te hebben onze gevoelens onder woorden te brengen en warme, ruimdenkende mensen te zoeken met wie we in veiligheid eindelijk onze onzekerheden kunnen toegeven - en onszelf door hun ogen kunnen leren zien: eerlijker, zonder oordeel en met meer compassie. Door hun liefde kunnen we hopen een manier te vinden om te uiten wat we verlangen en wie we zijn, met iets minder angst.
Zelfs stoppen met de verwachting dat seks eenvoudig voor ons zou kunnen zijn, is een enorme vooruitgang. Eenvoudig erkennen dat we moeilijkheden hebben op dit gebied is het begin van vooruitgang en bevrijding.
Om te meten hoeveel schaamte we in ons dragen, kunnen we onszelf enkele scherpe vragen stellen waarop misschien geen prettige antwoorden volgen:
Hoe voel je je over je eigen lichaam?
Hoeveel medelijden voel je met degene die seks met jou heeft?
Zou iemand jou seksueel echt kunnen kennen en je daarna nog leuk kunnen vinden?
Wij, de beschaamden, verdienen het om seks opnieuw te ontdekken. Niet als een zone van schuld en angst, maar als een intens vervullende, onschuldige en in wezen ‘leuke’ activiteit, iets wat we echt verdienen te ervaren - net zoals we, ondanks vroege aanwijzingen van het tegendeel, echt verdienen om te bestaan.

Toch is het cruciaal te herkennen wanneer een voorkeur een onverwerkte wond maskeert. Als iemand alleen veiligheid voelt in dynamieken die op vroegere pijn lijken, kan dit wijzen op noodzakelijk helingswerk. Voorkeuren kunnen emoties gezond verwerken, maar mogen geen vervanging zijn voor echte genezing.

Fantasie versus realiteit: waarom ze niet hetzelfde zijn
Een groot misverstand is dat seksuele fantasieën weerspiegelen wat iemand in het echt wil. Dat klopt niet altijd. Iemand kan non-consentfantasieën hebben, maar dit absoluut níet in werkelijkheid ervaren. Een ander geniet van onderwerping in de slaapkamer, maar is dagelijks onafhankelijk en assertief. Ons brein werkt nu eenmaal in compartimenten. Het onderscheid tussen gedachte en echt verlangen begrijpen, is cruciaal. Niet elke fantasie hoeft te worden uitgeleefd; sommige horen thuis in de verbeelding, en dat is prima. Gevaar ontstaat pas wanneer mensen impulsen blind volgen zonder zich af te vragen waar ze vandaan komen.

De rol van schaamte en onderdrukking in vermijding
Soms is wat we een "voorkeur" noemen, eigenlijk een afleiding van iets diepers. Iemand kan extreme seksuele ervaringen opzoeken om hun ware seksuele geaardheid te vermijden. Of iemand met onverwerkt trauma gebruikt intense seks om emoties te verdoven. De voorkeur zelf is niet het probleem – de reden waarom we ernaar handelen, telt. Zijn het uitingen van plezier en zelfexpressie? Of een ontsnapping uit iets wat we niet onder ogen willen zien? Dat is de kernvraag.

Waarom open gesprekken nodig zijn
Voor een gezondere relatie met seksualiteit moeten we er zonder schaamte over kunnen praten. Oordelen doen verlangens niet verdwijnen; ze drijven ze alleen de schaduw in. Zonder zelfreflectie verliezen we de kans onszelf écht te begrijpen. Door ruimte te creëren voor eerlijke dialogen doorbreken we mythes, verminderen we schade en geven we mensen tools om hun seksualiteit veilig, consensueel en bewust te verkennen.

Als je hier nieuwsgierig naar bent of ermee worstelt: je bent niet alleen, en je bent niet kapot. Hoe meer we onszelf begrijpen, hoe betere keuzes we maken – in elk aspect van het leven.

Het kan tegenwoordig vreemd lijken om schaamte te voelen over je interesse in of gevoelens rondom seks. Als moderne, verlichte mensen zouden we allemaal extreem zelfverzekerd, goed aangepast en enthousiast moeten zijn over het onderwerp seks.

Maar niets is minder waar. Seksuele schaamte is, in werkelijkheid, voor velen van ons helemaal niet verdwenen. Het is vooral een psychologisch probleem en niet zomaar een politiek of religieus vraagstuk.

Onze capaciteit om ons seksueel op een zelfverzekerde en gelukkige manier uit te drukken, onze mogelijkheid om te zeggen wat we willen, er zonder schaamte om te vragen en snel situaties te verlaten waarin we ons niet vervuld of vernederd voelen, zijn allemaal enorme psychologische prestaties. Deze zaken zijn vaak alleen spontaan toegankelijk voor degenen die een zeer ondersteunende en emotioneel ontwikkelde vroege omgeving hebben ervaren. Om van nature onbezwaarde, seksueel gelukkige volwassenen te zijn, is nodig dat anderen, die zelf ontspannen waren, ons lang geleden hebben laten voelen dat we acceptabel zijn zoals we zijn. Dat onze lichamen en hun functies natuurlijke en goede dingen zijn. Dat we niet stout of zondig waren omdat we nieuwsgierigheid uitten over lichamelijk genot. Dat het oké was om af en toe rommel te maken - en dat het bijvoorbeeld meer dan goed was om op de leeftijd van twee jaar echt gefascineerd en verrukt te zijn door het vreemde en wonderlijke bestaan van je eigen achterwerk.

Seksuele verlangens behoren tot de meest persoonlijke en kwetsbare dingen die we ooit worden uitgenodigd om te uiten, en ze stellen ons bloot aan mogelijke enorme niveaus van spot en schaamte. Zoals pestkoppen van alle soorten altijd hebben geweten: als je iemand snel wilt vernietigen, maak hem dan beschaamd over zijn seksualiteit; hij zal nooit meer het zelfvertrouwen hebben om je uit te dagen.

Er zijn weinig dingen die meer wezenlijk ‘onszelf’ zijn dan ons verlangen naar seksuele verbinding, en daarom komen onzekerheden over hoe aardig we zijn, hoe verdienend we zijn of hoe legitiem het is dat we bestaan, vaak sterk naar voren in de slaapkamer. Ze vernietigen ons vermogen om eenvoudige en onbewogen seksuele wezens te zijn. Om het simpel te stellen: als er enig risico bestaat dat we ons slecht voelen over onszelf, voelen we ons - door een psychologische onvermijdelijkheid - slecht over onszelf en seks. Wat seksuele problemen worden genoemd - zoals impotentie, vaginisme, een gebrek aan verlangen, schadelijke verslavingen of een algemene angst voor intimiteit - zijn eerst en vooral problemen van zelfhaat. In de regel kun je niet tegelijkertijd jezelf haten en een geweldige tijd in bed hebben.

Het begin van het oplossen van het probleem van seksuele schaamte berust op een basisacceptatie dat het probleem bestaat en waarschijnlijk al veel schade heeft aangericht in ons leven. We moeten leren het te benoemen en te erkennen, tegen onszelf en een paar dierbaren te zeggen: ‘Ik voel verlammende schaamte rondom seks, en dat is oké.’ Een toewijding aan verandering is wat telt. Ondanks alle vrolijke suggesties die het tegendeel beweren, lopen veel van ons, vrouwen en mannen, op dit moment (net als tijdens het hoogtepunt van de Spaanse Inquisitie) rond met intense seksuele schaamte. Dit is niet omdat wat we seksueel willen op objectieve manier ‘slecht’ is (dat wil zeggen, opzettelijk schadelijk voor iemand anders), maar omdat onze geschiedenis ons heeft gepredisponeerd om zo negatief te voelen over onszelf.

Een belangrijk effect van seksuele schaamte is dat het ons tot zwijgen brengt. We zijn zo beschaamd dat we zelfs onze schaamte niet kunnen uitspreken. Het is daarom van groot belang om het lef te hebben onze gevoelens onder woorden te brengen en warme, ruimdenkende mensen te zoeken met wie we in veiligheid eindelijk onze onzekerheden kunnen toegeven - en onszelf door hun ogen kunnen leren zien: eerlijker, zonder oordeel en met meer compassie. Door hun liefde kunnen we hopen een manier te vinden om te uiten wat we verlangen en wie we zijn, met iets minder angst.

Zelfs stoppen met de verwachting dat seks eenvoudig voor ons zou kunnen zijn, is een enorme vooruitgang. Eenvoudig erkennen dat we moeilijkheden hebben op dit gebied is het begin van vooruitgang en bevrijding.

Om te meten hoeveel schaamte we in ons dragen, kunnen we onszelf enkele scherpe vragen stellen waarop misschien geen prettige antwoorden volgen:

Hoe voel je je over je eigen lichaam?
Hoeveel medelijden voel je met degene die seks met jou heeft?
Zou iemand jou seksueel echt kunnen kennen en je daarna nog leuk kunnen vinden?

Wij, de beschaamden, verdienen het om seks opnieuw te ontdekken. Niet als een zone van schuld en angst, maar als een intens vervullende, onschuldige en in wezen ‘leuke’ activiteit, iets wat we echt verdienen te ervaren - net zoals we, ondanks vroege aanwijzingen van het tegendeel, echt verdienen om te bestaan.

Seksualiteit is een onderwerp dat zelden meer controverse, schaamte en verwarring oproept. Dit levert een probleem op wanneer mensen willen begrijpen wat gezonde seksualiteit is, want eerlijk gezegd is bijna niemand het erover eens. De reden hiervoor is dat meningen volledig worden gevormd door persoonlijke ervaringen, familie, maatschappij, religie en cultuur. Wat moeten we dan doen als we gezonde seks en een gezonde seksualiteit willen ervaren?

Wat is gezonde seksualiteit? Voordat we verder gaan, moet duidelijk worden dat er een enorm verschil is tussen genezen en gezond. Gezond impliceert dat je volledig genezen bent. Gezond impliceert dat de manier waarop je seksualiteit is, en seks als daad zelf, een weerspiegeling is van je perfect geïntegreerde zelf in je staat van uiterste menselijk potentieel. Als we gezonde seksualiteit willen beoefenen, moeten we het idee van 'gezond' loslaten, omdat we anders de noodzakelijke stappen overslaan om überhaupt tot gezonde seksualiteit te komen. Bovendien zullen we niet ontdekken wat ware gezonde seksualiteit voor ons specifiek is als we het van buitenaf benaderen, in plaats van van binnenuit.

De benadering van buitenaf, die we zo vaak toepassen, is om te beslissen hoe gezonde seks eruit moet zien en dat dan als een onzichtbare lat boven ons hoofd te houden, die we moeten bereiken. Dit is niet de juiste manier. De juiste manier is van binnenuit. De benadering van binnenuit houdt in dat we stappen blijven zetten in de richting van wat meer en meer helend aanvoelt, totdat we uiteindelijk uitkomen bij hoe gezonde seksualiteit eruitziet voor ons als unieke individuen. Het is het verschil tussen het hebben van een bestemming in gedachten en daarheen gaan, en het aangaan van een reis in de richting van wat meer en meer helend is om uiteindelijk op de bestemming aan te komen. We moeten ons richten op helende seksualiteit, niet op gezonde seksualiteit.

Op energetisch niveau gebeurt er tijdens helende seks dat de energie in het lichaam versnelt. Er wordt een uitlijning ervaren tussen je fysieke belichaming en je niet-fysieke belichaming. Dit zorgt ervoor dat al je chakra's en meridianen zich openen, zodat energie vrij kan stromen en je in een staat van integratie komt tussen je fysieke en niet-fysieke belichaming, evenals in uitlijning en integratie met een ander persoon.

Het leven op aarde, de schepping zelf, is het bijproduct van seksuele energie. Seksuele energie en conceptie/creatie gaan hand in hand, en je kunt dus zien dat seks en orgasme een element van extreme creatieve kracht bevatten. De vraag is: wat ben je aan het conceperen of creëren? Wat betekent dit voor degenen onder jullie die de kunst van manifestatie beoefenen? Het betekent dat seksuele focus een van de krachtigste hulpmiddelen is voor manifestatie die je je maar kunt voorstellen. Wat je focust, vooral op een gevoelsmatig niveau, is wat je probeert te conceperen en manifesteren in je realiteit. In het moment van een orgasme wordt die energie geaccumuleerd en vervolgens in een enorme uitbarsting uitgezonden. We kunnen die uitbarsting, die supernova van creatieve energie, gebruiken om bij te dragen aan wat we ook maar proberen te creëren.

Zelfs de meest fundamentele mensen op de planeet, de mensen die niet echt bezig zijn met zelfbewustzijn, de mensen die geen idee hebben van al deze energetische componenten van het leven, hebben een onderbewust begrip van de creatieve energie achter seksualiteit. Dit is een reden dat ze fetisjen hebben. Mensen hebben fetisjen omdat er een onmiskenbare link is tussen verlangen en seks. De realiteit is echter dat de meeste mensen geen bewustzijn hebben van wat hun werkelijke verlangens zijn. Ze voelen een verlangen naar iets of iemand, maar ze hebben geen bewust begrip van waarom ze dat specifieke verlangen hebben. Dit is wat er schuilgaat achter alle fetisjen. Het geheim achter alle fetisjen is dat er achter elk daarvan iets schuilt dat de persoon wil ervaren, meestal een emotionele toestand waarvan ze het gevoel hebben dat ze die volledig missen en wanhopig willen ervaren, oftewel creëren. Maar die persoon gelooft dat ze het niet direct kunnen creëren. Dus de fetisjen van mensen suggereren wat ze nodig hebben om te genezen; ze suggereren wat helend is voor iemand.

Genezing vindt plaats in een progressie. Het is duidelijk dat iemand die echt heel boos is, niet in een perfecte staat verkeert van wat we 'gezondheid' zouden noemen, omdat ze niet in een staat van vreugde of innerlijke rust zijn. Dat is vanzelfsprekend. Maar als iemand zich machteloos voelt, is boos worden absoluut noodzakelijk voor hun genezingsprogressie. Sterker nog, ze kunnen niet in een staat van vreugde en innerlijke rust komen zonder eerst die woede te kunnen voelen.

Mensen die met mij willen praten over wat gezonde seks is, hebben steevast vragen over het concept van machtsdynamiek in seks, zoals dominantie en onderwerping. Voor de context van deze video, zodat je begrijpt wat ik bedoel met gezonde seksualiteit, zal ik vaak dit voorbeeld van seksuele dominantie gebruiken. Laten we eens kijken naar dominantieseks. Bij dominantieseks heeft de ene persoon, die zich machteloos heeft gevoeld, nodig om in de controlerende positie te zijn tijdens de seksuele interactie om te genezen. Het is vrij gewoon dat de andere persoon, die vrijwillig deelneemt aan die machtsdynamiekseks en een onderdanige rol aanneemt, in hun leven oververantwoordelijk is geweest voor zaken. Ze hebben die druk voor leiderschap en die druk om dingen alleen te doen in zo'n extreme mate gevoeld dat het een ongelooflijke opluchting en een diepe behoefte is om al die verantwoordelijkheid los te laten en zich in feite te laten commanderen. Als deze twee mensen zich bewust bezighouden met deze handeling van seksuele dominantie en onderwerping, is het feitelijk helend voor beide betrokken partijen.

Nu, natuurlijk, degenen onder ons die zich in de spirituele wereld begeven, begrijpen dat machtsdynamiek voortkomt uit het ego. Het komt voort uit een behoefte om dingen te controleren. Dus hebben we van het ego een vijand gemaakt, en daarom kijken we naar iemand die dominantieseks nodig heeft met een houding als "dat is niet helemaal gezond", op dezelfde manier als we kijken naar iemand die boos is, in termen van "dat is niet helemaal gezond". We moeten ernaar kijken in termen van: "het is feitelijk helend, en daarom is het een noodzakelijke stap voor hen om dat te doen om in een gezonde ruimte te komen." Er is niets mis of minderwaardig aan die toestand, net zomin als er iets mis is met kruipen voordat je loopt en rent. Dus nu je dat punt begrijpt, laten we eens diep ingaan op waar helende seksualiteit over gaat.

Kenmerken van Helende Seksualiteit
Helende seksualiteit is precies dat: het is helend. Genezen betekent het tegenovergestelde ervaren. Helende seksualiteit zal ons dichter bij integratie brengen, in plaats van verder ervan af. Dit is waarom helende seksualiteit er voor alle mensen en alle koppels anders uit zal zien. Voor de één kan het helend zijn om zich bezig te houden met dominantieseks; voor de ander kan het juist het tegenovergestelde van helend zijn. We moeten ons dus heel bewust worden van of iets helend voor ons is, of schadelijk voor ons specifiek. Om te ontdekken wat helend is, moeten we het idee loslaten van welk type seks we zouden moeten hebben. Deze standaarden gaan over een idee van goed en kwaad. Het punt is dat onze beoordeling van goed en kwaad zo zelden voortkomt uit een authentieke beoordeling van wat persoonlijk juist en in lijn is voor ons. Maar al te vaak – ik heb het over de meeste tijd – is het het resultaat van beïnvloeding door allerlei mensen, omstandigheden, culturen, alles, om tot een idee van goed en kwaad te komen dat niet eens authentiek is voor ons. Onze seksualiteit moet voortkomen uit ons unieke, authentieke wezen.

Helende seksualiteit moet helend zijn voor alle betrokken partijen, inclusief onze eigen interne delen. Dit betekent dat helende seksualiteit consensueel is. Veel van de fetisjen die geen consensuele seks inhouden, hebben helemaal niets met seks te maken. Ze hebben te maken met iets anders. Een voorbeeld: necrofilie gaat niet echt over seks hebben met doden. Het gaat over het ontbreken van de optie tot afwijzing op welke manier dan ook. Het gaat over volledig en totaal ontvangen worden. Voor deze persoon betekent genezing het vinden van situaties waarin ze volledig ontvangen kunnen worden zonder de mogelijkheid van afwijzing. Voor mensen die pedofielen zijn, een andere niet-consensuele interactie, gaat het niet om seks hebben met een kind. Het gaat over de desperate behoefte om opnieuw contact te maken met de eigen onschuld, de eigen essentie, en daarmee het eigen innerlijke kind. Dus wat men moet doen, is geen seks hebben met een kind, maar opnieuw contact maken met de eigen essentie, onschuld en het innerlijke kind. Helende seks mag ook niet het ene deel van onszelf genezen, terwijl het een ander deel van onszelf schaadt. Bijvoorbeeld: voor de ene persoon kunnen alle delen van zichzelf perfect in lijn zijn in een gegeven scenario met onbeschermde seks, wat betekent geen condooms of preventieve maatregelen. Voor een andere persoon wil een deel van hen misschien de intimiteit van geen condoom, terwijl een ander deel van hen volledig beschadigd kan worden door het risico van een ongeplande zwangerschap of een soa. En wanneer dat het geval is, is het eigenlijk nog steeds niet-consensuele seks. Het deel dat niet toestemt, is echter een deel in onszelf. Veilige en gezonde seks met betrekking tot zaken als bescherming tegen soa's en zwangerschap enzovoort, gaat echt over dit punt op de lijst.

We moeten accepteren dat we seksuele wezens zijn. Dit is natuurlijk voor ons. We kunnen dit niet blijven ontkennen. En seks omvat veel meer dan alleen de daad van seks zelf. Helende seksualiteit betekent het omarmen en genieten van je seksualiteit gedurende je leven, en het is een belangrijk onderdeel van je emotionele, mentale en fysieke gezondheid. Om je seksualiteit te genezen, moet je de aspecten in jou onder ogen zien die seksualiteit weerstaan en tegenwerken, evenals de trauma's die die weerstand hebben gecreëerd. We moeten ons bewust worden van alles in ons dat seksualiteit en seks tegenwerkt. De delen die geen plezier associëren met seks of intimiteit met seks, en we moeten die weerstand beginnen op te lossen. Je moet openstaan voor het idee dat voor de ene persoon de helende stap met betrekking tot seksualiteit, alles gezegd hebbende wat ik net heb gezegd, is om helemaal geen seks te hebben. Dit is vaak het geval wanneer mensen te vaak niet-consensuele seks hebben ervaren, vooral wanneer die seksualiteit in de kindertijd werd geïntroduceerd. Meestal is een van de zeer cruciale stappen die een persoon kan zetten op het genezingspad van seksualiteit, nee zeggen tegen seks. Dat gezegd hebbende, het is onderdeel van de genezing. Als we ons leven lang nee zeggen tegen seks en nee tegen seksualiteit, genezen we niet echt, toch? We blijven vastzitten!

We moeten accepteren dat er geen tegenstrijdigheid is tussen seksualiteit en spiritualiteit, seksualiteit en moraliteit, seksualiteit en goedheid. Zoveel van wat ons in de war brengt over seks, is dit concept van schaamte. En ik moet eerlijk zijn, spiritualiteit is een echte boosdoener voor dit idee. We hebben het idee dat seksualiteit, dat gerelateerd is aan lichamelijkheid en vleselijke lust en driften, op de een of andere manier niet in orde is. En het is echt tijd dat dit eindigt: dit idee dat we onze menselijkheid moeten overstijgen, onze lichamelijkheid moeten overstijgen, en spiritualiteit moeten gebruiken om te ontsnappen aan alles wat ons menselijk maakt. Dat betekent dat een deel van helende seksualiteit is om te beseffen dat het omarmen van seksualiteit en seks deel moet uitmaken van ons spirituele pad.

Experimenteren. Voor de meeste mensen gaat het pad van helende seks en seksualiteit niet alleen over het vinden van de delen in je die seksualiteit tegenwerken en daarmee omgaan. Het gaat ook over het doorbreken van de beperkingen van je beperkingen met betrekking tot seksualiteit en seksuele ervaringen. We moeten experimenteren op een manier die veilig aanvoelt voor ons, die in overeenstemming is met onze beste belangen. We moeten onderzoek doen en nieuwe dingen proberen, en ook energie steken in het ontdekken van wat seksueel plezier voor ons specifiek is. Je hoeft niet per se elk voedselproduct te proberen dat op de planeet beschikbaar is, maar je kunt jezelf niet echt als gezond beschouwen, gerelateerd aan voedsel, als je ooit maar één soort voedsel op de planeet hebt gegeten.

We moeten seks hebben met mensen vanuit een ruimte van liefde. Nu, wacht even, want ik weet dat toen ik die uitspraak net deed, het beeld dat je in je hoofd had, de liefdesscènes waren die je in een Hallmark-film ziet. Dat is niet wat ik zeg. Liefhebben is iets als een deel van jezelf nemen. Wanneer je dat doet, worden hun beste belangen onderdeel van jouw beste belangen, en dit is waar helende seksualiteit echt kan plaatsvinden. Als ik van een persoon houd, dan neem ik hun beste belangen als onderdeel van mijn beste belangen, en ben ik echt afgestemd op hun beste belangen. Misschien is hun beste belang op dit moment om in een onderdanige rol te zijn, daarom kan ik liefdevol zijn en ze compleet domineren. Maar als ik afgestemd ben op iemand anders, dan ben ik echt afgestemd op waar ze zijn en wat hun beste belangen zijn, en het is duidelijk, zonder mijn projectie, dat die persoon niet in een afgestemde staat zou zijn of dat het niet helend voor hen zou zijn om in een onderdanige rol te zijn, dan mag ik die persoon niet domineren. We moeten voldoende afgestemd zijn om de beste belangen van de ander te voelen, te zien, te horen en te begrijpen, in plaats van onze ideeën over wat hun beste belangen zouden moeten zijn op hen te projecteren. Wat, geloof me, draait om onze eigen voorkeur.

We moeten seksualiteit en de manier waarop we seks hebben, laten muteren, vloeien en evolueren gedurende ons leven en onze relaties. De seks die we hebben, zal compleet anders zijn wanneer we iemand voor het eerst ontmoeten versus op ons 10-jarig jubileum. Het zal compleet anders zijn. Onze seksualiteit en onze relaties tot seks, en de seks die we hebben, wanneer we een peuter zijn, zeg maar, en we nauwelijks opmerken: "Ik ben een meisje, dat is een jongen." Wanneer we zeven zijn en we onze eerste aantrekkingskracht beginnen te vormen tot het andere geslacht, of zelfs begrijpen wat onze genderpreferenties zijn in termen van seksualiteit, zal dat anders zijn dan toen we 17 waren. Wat weer anders zal zijn dan toen we 30 waren. Wat weer anders zal zijn dan toen we 60 waren. Als we ons inzetten voor helende seksualiteit, zal onze seksualiteit muteren en meer en meer in lijn komen met ons hoogste potentieel, evenals een meer oprechte uitdrukking van ons ware zelf. Onze seksualiteit zal steeds meer geïntegreerd worden. Helende seksualiteit is een vloeiend en flexibel proces en praktijk, geen starre. Seksuele energie is creatieve energie, daarom is seksualiteit en seks een creatief proces.

We moeten ons aansluiten bij de beweging om seks en seksualiteit niet langer taboe te maken. Het feit dat seksualiteit en seks een taboe-onderwerp zijn binnen de samenleving, is de schuld van zoveel ongezondheid rond seksualiteit in het algemeen. Daarnaast is de schaamte die we hebben rond naaktheid en het taboe van onze naaktheid hetzelfde: het keert ons tegen onszelf. Het maakt een deel van onszelf slecht en beschamend. Dit moet veranderen. Onze seksualiteit is zo verdeeld over seks, en dit is wereldwijd, dit is een menselijke conditie, en deze splitsing die we hebben met betrekking tot seksualiteit schaadt ons mentaal, het schaadt ons emotioneel, het schaadt ons fysiek, het schaadt onze relaties ook, op een grote manier.

Het punt is: seks verkoopt. Overal waar je kijkt, vooral in de media, is het seks, seks, seks, seks, seks, toch? Dus ons wordt verteld dat seks je waardevol maakt. Je moet sexy zijn, seks is belangrijk, seks is alles. Tegelijkertijd ontvangen we de exacte tegenovergestelde boodschap. We ontvangen de boodschap dat onze naakte lichamen volledig bedekt moeten worden omdat het een schande is als je ze op welke manier dan ook blootgeeft. Ons wordt verteld dat seksualiteit vies is, dat het zondig is, dat het je geen fatsoenlijke burger maakt, dat het je een paria maakt. Ik bedoel, de lijst gaat maar door. Wat moeten mensen in godsnaam doen met deze gemengde boodschap? Je zit in een dubbele binding, iedereen zit erin. Seks en seksualiteit mogen niet taboe zijn, net zomin als onze naakte lichamen. Dit taboe-karakter van seks leidt tot onderdrukking, en onderdrukking leidt tot onbewustzijn. Onbewuste seks is de echte oorzaak van seksuele schade op de planeet vandaag de dag. Je kunt energie niet onderdrukken; het zal zich opbouwen totdat de energie ergens heen moet, maar als die energie op een onbewuste manier naar boven komt, zal het tot pijn leiden. En laat me je vertellen, dit is precies waarom priesters hun misdienaars misbruiken. Wanneer je in een samenleving terechtkomt die seks onderdrukt, zal de seksuele deviatie in die cultuur je de stuipen op het lijf jagen.

Dus we moeten stoppen met het taboe maken van seks. Het feit dat we niet openlijk en bewust praten over seks en seksualiteit is de reden voor al deze verwarring, de reden voor al deze pijn en al deze ongezonde seksualiteit. En ik moet ook zeggen, het wegnemen van dit taboe van seksualiteit omvat ook het praten met onze kinderen als ouders en als opvoeders over seks. Zodat we een persoon niet tegen hun eigen lichaam keren, zodat we een persoon niet tegen hun eigen seksualiteit keren, en zodat we hen kunnen leren dit natuurlijke aspect van zichzelf te omarmen zonder in een plaats van pijn te komen, en zodat ze seksuele keuzes voor zichzelf kunnen maken die specifiek in lijn zijn met henzelf.

Seksualiteit is uiteindelijk een deel van wie je bent. Het is tijd om het te omarmen. Het is een half geleefd leven als je je leven lang je seksualiteit en seks in combinatie met die seksualiteit weerstaat. Dus de tijd is gekomen, het is tijd om seksualiteit te omarmen. En de manier om dit te doen is niet door jezelf op te houden aan een belachelijk rare standaard die iemand anders voor je heeft ingesteld, die je van buitenaf hebt proberen te achterhalen over hoe gezonde seksualiteit eruitziet. In plaats daarvan, zet je in voor helende seksualiteit.
